% =====================================================================
% §2.4 — La deviazione per grano: ampiezza e probabilità
%         (esempi gemelli: ex3a_range, ex3b_dephase)
%
% Pattern: aggancio all'identità di contenuto lasciata aperta da §2.3 ->
%          diff gemelli (la sola posizione dell'envelope) -> lettura della
%          figura doppia -> meccanismo (tendency mask + gate ortogonale) ->
%          quadrato 2x2 (envelope puro / Truax / gated / micromodulazione-
%          Vaggione) -> filo del pettine (asse contenuto, complementare al
%          tempo di §2.3) -> riproducibilità (andamento) -> ponte a §2.5
%
% TODO prima della submission:
%   [ ] linerange dei due listati: allineare alle righe reali di
%       ex3a_range.yml (blocco pointer) e ex3b_dephase.yml (pointer+dephase)
%   [ ] verificare nomi figure: ex3a_range_map.pdf, ex3b_dephase_map.pdf
%   [x] descrittori uditivi rimossi (sfocatura/crepitio); "vocale" -> "punto
%       del buffer"; "ronzio" -> "periodicità"; "partiture" -> "map"
%   [ ] numeri della micromodulazione (±1,5 dB; ±15°; ±5 ms; ±5% buffer):
%       aggiornare se i default_jitter cambiano (issue pitch in corso)
%   [ ] \cite{Truax1988} e \cite{Vaggione2002}: allineare alle chiavi reali
%   [ ] verificare \ref{fig:pointer-MAP} e \ref{sec:griglia} contro i label reali
%   [ ] figure*: verificare resa a doppia colonna con cim2026.sty
% =====================================================================

% ----------------------------------------------------------------
\subsection{La deviazione per grano: ampiezza e probabilità}\label{sec:deviazione}

A lettura ferma i grani restano identici nel contenuto, e con la
\texttt{distribution} (\S\ref{sec:griglia}) ne abbiamo irregolarizzato
soltanto il tempo. Sciogliere anche l'identità di contenuto è la
deviazione per grano: attorno alla traiettoria dichiarata, ogni
parametro può ricevere, grano per grano, uno scarto campionato in modo
indipendente. Il sistema fattorizza questa stocasticità in due
dimensioni a loro volta indipendenti --- \emph{quanto} un grano può
deviare, e \emph{con quale probabilità} la deviazione avviene --- ed
entrambe accettano envelope. I due esempi di questa sottosezione sono
gemelli costruiti per isolarle: stessa base (lettura congelata su un
punto del buffer, come al termine di \S\ref{sec:pointer}), stessa banda
massima di deviazione; differiscono per la sola posizione dell'envelope.

\begin{figure}[t]
    \includegraphics[width=\linewidth]{examples/ex3a_range/ex3a_range_map}\\
    \captionof{figure}{Gemello A (\texttt{ex3a\_range}), assi come in
    Fig.~\ref{fig:pointer-MAP}: cresce l'\emph{ampiezza} della
    deviazione --- cuneo a riempimento uniforme.}
    \label{fig:gemello-range}
\end{figure}

\begin{figure}[t]
    \includegraphics[width=\linewidth]{examples/ex3b_dephase/ex3b_dephase_map}\\
    \captionof{figure}{Gemello B (\texttt{ex3b\_dephase}): a banda
    massima invariata cresce la \emph{probabilità} --- la linea centrale
    persiste mentre la popolazione deviante si infittisce.}
    \label{fig:gemello-dephase}
\end{figure}


Nel primo gemello l'envelope sta sull'\emph{ampiezza}
(Listato~\ref{lst:mask-range}): il range di deviazione della posizione
di lettura cresce da zero al 35\% del buffer, e --- in assenza di altre
indicazioni --- la maschera è sempre attiva: ogni grano devia.

    \lstinputlisting[
      linerange={36-42},
      language=yaml,
      basicstyle=\ttfamily\fontsize{7.5}{9}\selectfont,
      label=lst:mask-range,
      caption={Gemello A: l'envelope sull'ampiezza della deviazione
      (\texttt{ex3a\_range.yml}).}
    ]{examples/ex3a_range/ex3a_range.yml}

Nel secondo gemello l'envelope si sposta sulla \emph{probabilità}
(Listato~\ref{lst:mask-dephase}): l'ampiezza resta fissa al massimo del
gemello A, e a crescere da 0 a 100\% è la chiave \texttt{dephase} --- la
frazione di grani a cui la deviazione si applica.

    \lstinputlisting[
      linerange={36-43},
      language=yaml,
      basicstyle=\ttfamily\fontsize{7.5}{9}\selectfont,
      label=lst:mask-dephase,
      caption={Gemello B: ampiezza fissa, envelope sulla probabilità
      (\texttt{ex3b\_dephase.yml}).}
    ]{examples/ex3b_dephase/ex3b_dephase.yml}

La Fig.~\ref{fig:gemello-range} e~\ref{fig:gemello-dephase} mostrano due
\emph{map}, e le due morfologie non potrebbero essere più diverse. Nella
prima un cuneo: la banda si allarga con continuità e i grani la riempiono
uniformemente, nessuno resta sulla linea di partenza perché tutti deviano
--- di poco all'inizio, di molto alla fine. Nella seconda la linea
centrale persiste per l'intera durata, mentre una popolazione crescente
la abbandona saltando da subito sull'intera banda: non un allargamento ma
una mistura, grani fedeli e grani devianti in proporzione componibile nel
tempo. Stessa base, stessa banda massima: ampiezza e probabilità della
deviazione sono assi indipendenti del controllo.

Il meccanismo sottostante è la \emph{tendency mask} della tradizione
granulare~\cite{Truax1988}: per ogni parametro una traiettoria centrale
nel tempo, un range di deviazione e un campionamento indipendente per
grano (uniforme, centro~$\pm$~range/2; una distribuzione gaussiana è
selezionabile in alternativa). Il gemello A è questo modello classico
all'opera. Il contributo proprio del sistema è il secondo asse: un gate
probabilistico, anch'esso componibile come envelope, che decide grano
per grano \emph{se} la maschera si applica --- il gemello B, dove non è
una banda a deformarsi ma una popolazione a migrare.

Il quadro delle combinazioni si lascia ora leggere per intero. Senza
range né dephase, il parametro segue l'envelope puro: è il regime
interamente scritto di \S\ref{sec:pointer}. Range senza dephase: la
maschera sempre attiva di Truax (gemello A). Range con dephase: la
maschera \emph{gated} (gemello B). Resta l'ultimo angolo, il dephase
senza alcun range: il gate apre allora deviazioni implicite di entità
minima definite dal sistema --- circa $\pm$1,5~dB sul volume,
$\pm$15\textdegree~sul pan, $\pm$5~ms sulla durata, $\pm$5\% del buffer
sulla lettura --- una micromodulazione che non disegna alcuna
traiettoria ma decorrela la massa grano per grano, nei termini della
\emph{décorrélation microtemporelle} di Vaggione~\cite{Vaggione2002}. È
la controparte, sull'asse del contenuto, di ciò che la
\texttt{distribution} opera sul tempo (\S\ref{sec:griglia}):
irregolarizzando ciò che ciascun grano legge, disperde la stessa
periodicità che il congelamento imponeva. Da sola, senza nulla di
scritto, basta a scioglierla (esempio audio \texttt{exA\_micro} nel
bundle).

Una parola sulla riproducibilità, qui che la stocasticità è in primo
piano: due rendering della stessa specifica producono grani diversi ma
la stessa morfologia. Ciò che lo YAML determina --- ed è lo YAML che il
bundle spedisce --- è l'andamento; la realizzazione allegata ne è
un'istanza.

Tutto ciò che precede riguarda però ancora una voce sola: un'unica
testina di lettura, un'unica griglia di emissione, una massa di grani
indifferenziata. Lo scostamento successivo moltiplica la voce.